\section{INTRODUCTION}
\label{sec:intro}

Imaging-based spatial transcriptomics assays such as Xenium \citep{janesick2023} measure thousands of transcripts per cell in situ, at single-cell resolution, over sections containing hundreds of thousands of cells. The appeal is obvious: a cell's expression can be read together with the identity and state of the cells around it, which is the raw material for asking how a microenvironment shapes the cells within it. The difficulty is that a cell's measured count vector is not a clean readout of that cell. It is at least three things superimposed. First, the cell's \emph{intrinsic} state: its type, its cell-cycle phase, its clone-associated expression, and the systemic signals that reach every cell of its type regardless of position. Second, its \emph{response} to the local microenvironment: the programmes that differ with who its neighbours are and what the surrounding tissue architecture looks like. Third, a \emph{measurement artefact} specific to in situ assays: transcripts that were detected in the vicinity of a segmentation boundary and assigned to the wrong cell. Segmentation of densely packed tissue is imperfect, and published estimates place the fraction of misassigned transcripts anywhere between one tenth and one half \needsource.

The second and third components are confounded at first order. A cell whose counts are contaminated by its neighbours acquires their markers; a cell that responds to its neighbours shifts its own programmes in a way their types predict. Where the responding genes are also markers of the neighbouring types, as for many ligands and receptors, or where a programme is shared across a region, the two produce the same signature, and neighbour composition predicts both. That misassigned transcripts produce apparent neighbourhood effects, including spurious cell--cell interaction signals, has been shown directly on in situ data \citep{mitchel2026celladmix}, and the contamination a cell receives grows with the local density of the donor type \citep{bilous2026split}. The two differ in form, one adds a neighbour's own profile, the other shifts the cell's own programmes, but from the cell-by-gene count matrix alone the split between them is fixed only by the parametric forms one is willing to assume, and a flexible decoder can mimic either. Any claim that a spatial effect on expression is real therefore depends, silently, on how much of the neighbour-predictable signal was attributed to contamination.

This paper makes that dependence explicit rather than hiding it. We propose DISCELL (\emph{disentangled cell states}), a generative model of per-cell counts with two latent variables and one fixed mixing operator. A latent $\vz_i$ carries the cell's intrinsic state and is, by architecture and by regularisation, kept free of information about the microenvironment given the cell's type. A latent $\vw_i$ carries the variation associated with a deterministic context descriptor $\vc_i$ that summarises who and what is nearby. A sparse, precomputed leakage kernel $\beta$ mixes a fraction $\kappa$ of each neighbour's clean expression profile into the cell's expected composition. The leak fraction $\kappa$ is not estimated. It is fixed, the model is refitted on a grid of $\kappa$ values, and every spatial effect is reported as a function of $\kappa$ with its envelope over random seeds, a sensitivity analysis in which $\kappa$ is the sensitivity parameter. An effect that is stable across the grid and across seeds cannot be explained by leakage of the modelled form; one that changes sign or vanishes is not separable even from that. The range across the grid is the deliverable, reported alongside an operating point. It is not a confidence interval, and it does not shrink with more cells.

Our modelling choices follow from three requirements that we motivate in \cref{sec:method} and \cref{app:rationale}. The decoder receives no type label, so that $\vz$ is forced to carry identity and cannot degenerate into a residual code. The context encoder queries the neighbourhood with the cell's type rather than with $\vz_i$ and reads only the neighbours' types, closing a channel through which attention would otherwise reconstruct the cell from its most similar neighbours and make the prior on $\vw$ ego-informed. And the invariance we enforce is \emph{conditional} on type: cell types cluster in space, so a marginal invariance between $\vz$ and neighbourhood composition deletes cell type from $\vz$ altogether, the trade-off known from demographic parity \citep{zhao2019inherent} and observed for conditional subspace VAEs \citep{slavutsky2025}.

\paragraph{Related work.}
Node-centric expression models \citep{fischer2023} regress a cell's expression on its neighbourhood directly and produce a niche descriptor without a per-cell latent; our context descriptor $\vc_i$ plays that role, and $\vw_i$ is what remains after conditioning on it. Several models share our design goal of separating intrinsic from spatially associated variation. SIMVI \citep{dong2025simvi} splits a cell's state into an intrinsic and a spatially induced part with an independence regulariser; MintFlow \citep{akbarnejad2025mintflow} conditions an intrinsic latent on cell type and a microenvironment latent on neighbour-type composition, and discourages the former from predicting the latter with type-specific discriminators; Celcomen \citep{megas2026celcomen} separates intra- from inter-cellular gene-interaction programmes with a generative graph neural network; and NicheCompass \citep{birk2025nichecompass} characterises niches with a graph neural network. None of them models transcript leakage, the confound this paper is about, and MintFlow names its susceptibility to segmentation errors as a limitation \todo{verify the NicheCompass characterisation; add scVIVA and further spatial VAEs}. Contamination-aware models for in situ data, notably resolVI \citep{ergen2025}, write a cell's counts as a mixture of its own decoded expression, its neighbours' expression decoded by the same network and weighted under a distance-kernel prior, and a slide-wide background, and estimate the mixing and neighbour weights per cell. Our leakage term is the same neighbour-mixing construction with the weights fixed, one leak fraction per section and no background term. We also adopt resolVI's stop-gradient on the neighbour path, which it uses for training stability, and describe the degeneracy it prevents (\cref{sec:batching}). resolVI separates the components by expecting true expression to be low-dimensional. A response to the neighbourhood is low-dimensional too, and the part of it that moves a cell toward its neighbours' profile has the same effect as leakage, so that expectation no longer separates the two on its own; we therefore sweep the leak fraction rather than estimate it \todo{support with the planted-leakage synthetic run once the synthetic appendix returns}. Variational disentanglement of condition-specific from shared representations \citep{klys2018,slavutsky2025} supplies the two-bound objective and the posterior-regularisation view of an adversarial invariance term \citep{ganchev2010} that we build on. We depart from them in three ways. The condition is not an observed label but the cell's niche, summarised by a deterministic context descriptor on which the prior of $\vw$ is conditioned. The invariance is conditional on cell type rather than marginal, because type and niche are dependent in tissue, as is standard in fair and domain-invariant representation learning \citep{madras2018,li2018ciddg,pogodin2023circe}. And our intrinsic path keeps the full decoder and draws $\vw$ from its prior instead of decoding from $\vz$ alone (\cref{app:pathb}). Our invariance target is a fixed data descriptor of the niche, so the model cannot shape it to make invariance trivial, unlike an independence penalty between two learned latents \citep{dong2025simvi}. That adversarial removal is not sufficient without an independent check is the observation of \citet{elazar2018}; our probe protocol is its application. \todo{extend related work: segmentation and contamination correction methods, further spatial VAEs and GNN spatial models}

\paragraph{Contributions.}
(i) A generative model that decomposes in situ counts into intrinsic state, microenvironmental response and neighbour leakage, with a multinomial likelihood in which the leakage mixture is exact under Poisson sources whose leaked-in rate scales with the receiving cell's own rate rather than the donors' (\cref{sec:generative}). (ii) An inference scheme whose architecture closes information channels that a regulariser cannot reach on its own: the intrinsic encoder never sees the context, the attention context is queried by type, and neighbour quantities are frozen, with an exact two-hop batching scheme that makes the frozen path affordable (\cref{sec:inference,sec:batching}). (iii) A conditional-invariance regulariser, adversarial in its operating configuration with a closed-form penalty as the cheaper first attempt, conditional on type as in the type-specific discriminators of \citet{akbarnejad2025mintflow} but targeting an image-derived niche descriptor alongside neighbour composition, together with a held-out probe protocol that grades either independently of the training objective (\cref{sec:invariance}). (iv) The leakage sweep as a sensitivity analysis: spatial effects are reported as trajectories over $\kappa$ with seed envelopes, alongside an operating point (\cref{sec:sweep}). \todo{(v) the empirical evaluation, to be added}
